\chapter{Conclusiones}\label{cap:conclusiones}

Este Trabajo Fin de Máster ha abordado la detección binaria de música de origen humano frente a música generada por IA mediante la comparación de dos enfoques bajo un mismo diseño experimental: un baseline basado en MFCC, \textit{StandardScaler} y SVM con núcleo RBF, y un enfoque basado en embeddings de MERT-v1-95M congelado, \textit{StandardScaler} y SVM lineal. El objetivo general se ha cumplido al evaluar ambos modelos con las mismas particiones de AIME y mediante medidas predictivas y operacionales, lo que ha permitido realizar una comparación acotada y seleccionar el modelo integrado en la prueba de concepto web.

Los objetivos específicos también quedan cubiertos. La revisión del estado del arte permitió delimitar la detección de canciones sintéticas como un problema distinto de tareas próximas y situar la comparación entre características acústicas clásicas y representaciones profundas preentrenadas. La auditoría de AIME permitió verificar su viabilidad para el experimento, con las limitaciones identificadas durante el análisis, y construir un subconjunto equilibrado de 1\,000 ejemplos, con 500 audios humanos de MTG-Jamendo y 500 generados por IA. Asimismo, se definió un protocolo reproducible de particionado, preprocesamiento y evaluación: las particiones \textit{train}, \textit{val} y \textit{test} se generaron por grupos de \textit{description}, con semilla fija y sin solapamiento entre grupos, y ambos enfoques se evaluaron sobre el mismo \textit{test} de 150 ejemplos.

Se implementó y evaluó el baseline MFCC + SVM, y se seleccionó, adaptó y evaluó MERT como encoder congelado para obtener embeddings musicales. En la evaluación final, MFCC + SVM obtuvo una \textit{balanced accuracy} de 0,8333, un F1 IA de 0,8299 y un ROC-AUC de 0,9086, frente a 0,8200, 0,8212 y 0,8985 de MERT + SVM, respectivamente. El baseline alcanzó 125 aciertos sobre 150, frente a los 123 de MERT, mientras que MERT obtuvo un \textit{recall} IA ligeramente superior, de 0,8267 frente a 0,8133.

Estos resultados muestran que ambos enfoques ofrecen resultados próximos dentro del \textit{test} de AIME. MFCC presenta una ligera ventaja en la mayoría de las medidas finales consideradas, mientras que MERT conserva una ligera ventaja en \textit{recall} de la clase IA. Dada la magnitud reducida de las diferencias y al no haberse realizado una prueba inferencial, estos resultados no permiten afirmar una superioridad estadísticamente significativa de MFCC sobre MERT. MERT fue viable en CPU, obtuvo mejores métricas en validación y mantuvo un rendimiento muy cercano al baseline en la evaluación final.

La comparación permitió cumplir el objetivo de seleccionar un modelo para la prueba de concepto. MFCC + \textit{StandardScaler} + SVM RBF se eligió para VeriSon como una decisión de ingeniería, apoyada tanto en sus resultados de \textit{test} ligeramente favorables como en su menor complejidad operacional. Su pipeline es más simple, no requiere ejecutar un encoder profundo durante la inferencia, incorpora menos componentes y dependencias, y facilita la integración y el despliegue. La selección responde, por tanto, al compromiso entre rendimiento predictivo y viabilidad operacional considerado más adecuado para el alcance de la aplicación desarrollada.

El objetivo software también se ha completado mediante el desarrollo y despliegue de VeriSon. La aplicación integra el modelo seleccionado en un backend FastAPI y una interfaz construida con React, Vite y TypeScript. El backend se empaqueta mediante Docker y se despliega en Northflank, mientras que el frontend se despliega en Vercel. Asimismo, se ha implementado integración continua con GitHub Actions para comprobar las pruebas del backend, las verificaciones del frontend y la construcción de la imagen Docker. De este modo, el modelo seleccionado se ha integrado en una aplicación web funcional, completando la vertiente de desarrollo software planteada en los objetivos del TFM.

Finalmente, las decisiones metodológicas, las configuraciones experimentales y los resultados obtenidos se han documentado y conservado en el repositorio junto con los artefactos necesarios para reconstruir el procedimiento seguido. Con ello se aporta trazabilidad al desarrollo y se cumple el objetivo de documentación y reproducibilidad del trabajo.

\section{Limitaciones y trabajo futuro}

Las conclusiones deben interpretarse dentro del alcance experimental definido. La evaluación se ha realizado sobre un único dataset y sobre un subconjunto de 1\,000 ejemplos. La clase humana procede exclusivamente de MTG-Jamendo y pueden persistir atajos asociados a la procedencia, el códec, la frecuencia de muestreo, la producción u otras diferencias técnicas entre fuentes. Además, \textit{description} es una agrupación semántica útil para el particionado, pero no un identificador verificable de canción original.

La evidencia experimental corresponde a clips de 10~s. Aunque VeriSon segmenta los archivos y agrega las puntuaciones de decisión de sus fragmentos para emitir una estimación a nivel de archivo, dicha agregación no equivale a una validación experimental sobre canciones completas.

Como continuación prioritaria, sería necesario evaluar ambos enfoques con datasets externos y con generadores que no hayan estado representados durante el desarrollo experimental, para estudiar su capacidad de generalización fuera del escenario cerrado utilizado. También conviene diseñar una evaluación específica a nivel de canción completa, analizando cómo afectan la segmentación y la agregación de decisiones de fragmentos al resultado final. Estas dos líneas permitirían contrastar de forma más rigurosa la utilidad del detector en condiciones más cercanas a un uso real.
